\label{ch:modular}

\section{Hexad--octad inclusion and associated functional realizations}

Let us fix a set \(H\) of six points and a set \(O\) of eight
points with a distinguished inclusion \(j:H\hookrightarrow O\). In the
Witt residence, \(H\) is a hexad of a dodecad and \(O\) its lifted octad;
in this chapter only the already fixed inclusion is used, not the
terminal value of Barbero--Immirzi. We write
\begin{equation}
\label{eq:modular-pairs-H}
\binom H2:=\{P\subset H:|P|=2\},
\qquad \left|\binom H2\right|=\binom62=15,
\end{equation}
and we define two incidence spaces with distinct types:
\begin{equation}
\label{eq:modular-F6-F8}
\mathcal F_6:=H\times\binom H2,
\qquad
\mathcal F_8:=O\times\binom H2.
\end{equation}
The inclusion of incidences is
\begin{equation}
\label{eq:modular-iota68}
\iota_{6,8}:\mathcal F_6\hookrightarrow\mathcal F_8,
\qquad
(h,P)\longmapsto(j(h),P).
\end{equation}
It is injective because \(j\) is injective and preserves the marked pair \(P\).

\begin{theorem}[Incidence and defect counts oriented]
\label{thm:modular-incidence-90-120}
The two spaces of \cref{eq:modular-F6-F8} satisfy
\begin{equation}
\label{eq:modular-counts-90-120}
|\mathcal F_6|=6\binom62=90,
\qquad
|\mathcal F_8|=8\binom62=120,
\qquad
|\mathcal F_8\setminus\iota_{6,8}(\mathcal F_6)|=30.
\end{equation}
Normalization by the twenty-four coordinates of the
dodecad--twenty-four pair and by the fifteen pairs of \(H\) produces
\begin{equation}
\label{eq:modular-incidence-defect}
\boxed{
\delta_{6,8}^{\rm inc}
=\frac{|\mathcal F_6|-|\mathcal F_8|}
{24\binom62}=-\frac1{12}.}
\end{equation}
The same fraction is obtained through the reciprocal functional associated
with the incorporated memory,
\begin{equation}
\label{eq:modular-reciprocal-defect}
\boxed{
\delta_{6,8}^{\rm rec}
=\frac{30}{120}-\frac{30}{90}=-\frac1{12}.}
\end{equation}
They are two exact realizations of the same oriented defect, but have
distinct domains and are not identified with \(\zeta(-1)\) without an
additional intertwiner.
\end{theorem}

\begin{proof}
The multiplicative principle applied to
\cref{eq:modular-F6-F8} gives \(6\cdot15=90\) and \(8\cdot15=120\).
The image of \(\iota_{6,8}\) has ninety elements, so the
complement has thirty. Substitution into
\cref{eq:modular-incidence-defect} gives
\(-30/(24\cdot15)=-1/12\), while
\cref{eq:modular-reciprocal-defect} is
\(1/4-1/3=-1/12\).
\end{proof}

The hexad--octad transition admits two functional realizations that
are not deduced from one another. The algebraic diagram, formulated on the spaces
that justify the counts, is
\begin{equation}
\label{eq:modular-68-two-readers}
\begin{aligned}
(H\xhookrightarrow{\,j\,}O)
&\longmapsto
\left(|\mathcal F_6|,|\mathcal F_8|\right)=(90,120)
\longmapsto\Gamma_{6\to8}=\gammaH,\\
(h_6,o_8,t)=(6,8,3)
&\longmapsto-\frac{o_8-h_6}{t}=-\frac23,
\end{aligned}
\end{equation}
where \(-2/3\) is the coefficient of the third coordinate of row \(13\) of the CKM constructor. The first realization transforms the incidence into area and modular memory; the second transforms the same transition into angular-mixing parameters. \(\boldsymbol\theta_{\rm CKM}=f(\gammaH)\) is not postulated: both outputs preserve a common antecedent.

The morphism that determines the sectors \(90/120\) is
\(\iota_{6,8}\). The quantity \(\gammaH\) has already been determined internally
by the hexad--octad construction and is reused here without reconstructing it from
its decimal representation. The
operatorial consequences begin only after fixing the elementary
scale \(a_0\), which is defined below without using entropy or the
terminal product \(\gammaH a_0\).

\section{Correspondence between the Nonadic Interior and the Boundary Degrees of Freedom}

Let \(\mathbb T_3=\{0,1,2\}\). Every \(x\in\mathbb Z/9\mathbb Z\) has
a unique digit representation
\(x=x_0+3x_1\), with \(x_i\in\mathbb T_3\), and every
\(y\in\mathbb Z/27\mathbb Z\) a unique representation
\(y=y_0+3y_1+9y_2\). Therefore
\begin{equation}
\label{eq:modular-bulk-boundary-729}
\mathcal B=(\mathbb Z/9\mathbb Z)^3,
\qquad
\mathcal B_\partial=(\mathbb Z/27\mathbb Z)^2,
\qquad
|\mathcal B|=|\mathcal B_\partial|=3^6=729.
\end{equation}
Equality of cardinalities is accompanied by the explicit
map. If \(x_j=r_{2j-1}+3r_{2j}\) for \(j=1,2,3\), we define
\begin{equation}
\label{eq:modular-beta729}
\beta_{729}(x_1,x_2,x_3)
=\bigl(r_1+3r_2+9r_3,\;r_4+3r_5+9r_6\bigr).
\end{equation}

\begin{proposition}[Genealogical bijection \(729\leftrightarrow729\)]
\label{prop:modular-beta729}
The map \(\beta_{729}:\mathcal B\to\mathcal B_\partial\) is bijective and
preserves the ordered six-trit word. In general, it is not an
isomorphism of additive groups: regrouping the digits changes where the
carry is recorded.
\end{proposition}

\begin{proof}
Extracting the six digits from three nonadic coordinates and grouping them into
two triples is a permutation of the coordinates of \(\mathbb T_3^6\). The
positional decomposition is unique on both sides, so the
operation has an inverse. The additive caveat is verified, for example,
by adding words that generate a carry between \(r_3\) and \(r_4\): that
boundary separates the two coordinates of order twenty-seven, but crosses
two distinct coordinates in the nonadic packing.
\end{proof}

On the boundary torus \(\mathcal B_\partial\), we take the block of vertices
\begin{equation}
\label{eq:modular-block-A}
A=\{0,1,\ldots,8\}\times\{0,1,\ldots,8\}.
\end{equation}
Its oriented links crossing the boundary are separated by direction:
\begin{align}
\partial A_{90}
&=\{((-1,j),(0,j)),((8,j),(9,j)):0\leq j\leq8\},
\label{eq:modular-boundary90}\\
\partial A_{120}
&=\{((i,-1),(i,0)),((i,8),(i,9)):0\leq i\leq8\},
\label{eq:modular-boundary120}
\end{align}
where all coordinates are read modulo twenty-seven. Thus
\begin{equation}
\label{eq:modular-boundary-counts}
\partial A=\partial A_{90}\sqcup\partial A_{120},
\qquad
|\partial A_{90}|=|\partial A_{120}|=18.
\end{equation}
The subscripts \(90,120\) label the channel provenance; they do not say that
the boundary has ninety or one hundred and twenty links.

\begin{theorem}[Law areal of the cut triadic]
\label{thm:modular-triadic-cut}
In the cubic graph \(N\times N\times N\) with unit capacity, every
cut between two opposite faces has capacity at least \(N^2\), and a
transverse layer attains that bound. For \(N=9\),
\begin{equation}
\label{eq:modular-mincut-entropy}
\mathcal A_{\min}=81,
\qquad
S_{\rm tri}=81\log3.
\end{equation}
\end{theorem}

\begin{proof}
There are \(N^2\) lines of edges parallel to the normal direction, pairwise
disjoint. Every cut must intersect each one, so its capacity
is at least \(N^2\). A single transverse layer contains exactly
\(N^2\) edges and attains the minimum. If each cut link carries a
maximally mixed trit, its entropy is \(\log3\); additivity over the
eighty-one links gives \cref{eq:modular-mincut-entropy}.
\end{proof}

The entropy \(S_{\rm tri}\) measures the capacity of the cubic cut. The
modular entropy of the thirty-six links of
\cref{eq:modular-boundary-counts} measures another state and must not be equated with
it.

\section{Collective incidence entangler between interior and boundary}

In the sum
\(\ell^2(\mathcal F_6)\oplus\ell^2(\mathcal F_8)\), we define
\begin{equation}
\label{eq:modular-u90-u120}
u_{90}=\frac1{\sqrt{90}}\sum_{f\in\mathcal F_6}(\delta_f,0),
\qquad
u_{120}=\frac1{\sqrt{120}}\sum_{f\in\mathcal F_8}(0,\delta_f).
\end{equation}
In
\(\ell^2(\partial A_{90})\oplus\ell^2(\partial A_{120})\), let
\begin{equation}
\label{eq:modular-v90-v120}
v_{90}=\frac1{\sqrt{18}}\sum_{e\in\partial A_{90}}(\delta_e,0),
\qquad
v_{120}=\frac1{\sqrt{18}}\sum_{e\in\partial A_{120}}(0,\delta_e).
\end{equation}
The pairs \((u_{90},u_{120})\) and \((v_{90},v_{120})\) are orthonormal.
Therefore there exists a unique isometry between the collective subspaces
that satisfies
\begin{equation}
\label{eq:modular-J68}
\mathcal J_{6\to8}u_{90}=v_{90},
\qquad
\mathcal J_{6\to8}u_{120}=v_{120}.
\end{equation}
Define
\begin{align}
D_{\rm inc}
&=90|u_{90}\rangle\langle u_{90}|
+120|u_{120}\rangle\langle u_{120}|,
\label{eq:modular-Dinc}\\
D_{\rm area}
&=90|v_{90}\rangle\langle v_{90}|
+120|v_{120}\rangle\langle v_{120}|.
\label{eq:modular-Darea}
\end{align}

\begin{theorem}[Exact transport of the two labels]
\label{thm:modular-J68}
The map \(\mathcal J_{6\to8}\) is unitary between the two collective
planes and intertwines the incidence and area operators:
\begin{equation}
\label{eq:modular-J-intertwining}
\boxed{
\mathcal J_{6\to8}D_{\rm inc}
=D_{\rm area}\mathcal J_{6\to8}.}
\end{equation}
Therefore, every polynomial combination of the labels \(90,120\), in
particular the primitive combination \(12:-1\), is transported with the
same coefficients.
\end{theorem}

\begin{proof}
Orthonormality shows that \(\mathcal J_{6\to8}\) maps one
orthonormal basis to another. In \(u_{90}\), both sides of
\cref{eq:modular-J-intertwining} equal \(90v_{90}\); in \(u_{120}\),
both equal \(120v_{120}\). Equality on a basis proves the identity.
The functional calculus, in its polynomial form, preserves the intertwining.
\end{proof}

This map does not purport to inject ninety individual incidences into
eighteen links: it transports exactly the two uniform modes and their
spectral labels. Extending it to non-uniform fluctuations would require
a new incidence map; the modular operator defined here only uses
the collective subspace that has just been constructed.

\section{Determination of the Elementary Scale of the Spectrum of the Area Operator}

Let \(C^*\) the angular coordinate conjugate, expresada in grados, and fijemos
\begin{equation}
\label{eq:modular-theta-clock-a0}
\theta_{\rm clk}=\frac{C^*}{6}\frac{\pi}{180},
\qquad
\boxed{
a_0=\frac{1+2\cos(10^\circ)}{3}\,\theta_{\rm clk}.}
\end{equation}
The factor preceding \(\theta_{\rm clk}\) is not a decimal
correction. It is the normalized real character of the phase triplet
\begin{equation}
\label{eq:modular-c10-character}
\frac13\Tr\operatorname{diag}
\left(1,\ee^{\ii\pi/18},\ee^{-\ii\pi/18}\right)
=\frac{1+2\cos(10^\circ)}3.
\end{equation}
Thus, \(a_0\) is the spectral mean of the angular structure on the oriented
trit. It is positive, is fixed before the area operator, and is not
selected from entropy. The subsequent identification
\(\gammaH=\Gamma_{6\to8}\) realizes the HMT output in the Ashtekar--Barbero connection
\cite{Barbero,Immirzi}; it does not reopen its derivation.

\begin{remark}[Alcance of the simbolo \(a_0\)]
In this chapter, \(a_0\equiv a_0^{\rm area}\) denotes exclusively the
normalization of \cref{eq:modular-theta-clock-a0}. It is neither the Bohr radius
nor the initial cosmological scale factor, although those objects are traditionally written
with the same symbol. No identity in this chapter
allows one to substitute one for the other.
\end{remark}
